\begin{frame}[fragile]
  \frametitle{Conclusiones}
  
  \begin{block}{Contribuciones principales}
    \begin{itemize}
      \item Demostración exitosa de PINNs en reconstrucción meteorológica a escala local
      \item Metodología reproducible y escalable para regiones con datos dispersos
      \item Validación de coherencia física ($\mathcal{R}_{NS} = 0,01273$) en modelos híbridos
    \end{itemize}
  \end{block}
  
  \begin{block}{Logro de objetivos}
    [Org] Organización y preprocesamiento: 16.368 registros de 18 estaciones\\
    [Diseño] Diseño e implementación: PINN con restricciones N-S integradas\\
    [Eval] Evaluación: 9 experimentos con métricas reproducibles
  \end{block}
  
  \note[item]{Tiempo: 2 minutos}
\end{frame}

\begin{frame}[fragile]
  \frametitle{Trabajo Futuro}
  
  \begin{block}{Extensiones inmediatas}
    \begin{itemize}
      \item Predicción temporal: transición de interpolación a pronóstico ($t + \Delta t$)
      \item Incorporación de variables adicionales: humedad relativa, radiación solar
      \item Análisis de sensibilidad a resolución de malla ($R$ más fina)
    \end{itemize}
  \end{block}
  
  \begin{block}{Aplicaciones emergentes}
    \begin{itemize}
      \item Extensión a otras regiones oceanográficas (Pacífico colombiano, cuencas andinas)
      \item Acoplamiento con modelos de dispersión de contaminantes
      \item Integración en sistemas operativos de monitoreo ambiental
    \end{itemize}
  \end{block}
  
  \note[item]{Tiempo: 1,5 minutos}
\end{frame}

\begin{frame}[plain,c]
  \centering
  \huge{Gracias}
  \note[item]{Cierre (0,5 min)}
\end{frame}
